% LuaLaTeX document - compile with: luatex-pdf dialog_design_principles.tex
\documentclass[10pt,a4paper,twocolumn]{ltjsarticle}

% ファイル生成日時（JST）
\newcommand{\generatedDate}{2026-01-12}
\newcommand{\generatedTime}{18:30}

% 基本パッケージ
\usepackage{geometry}
\geometry{left=15mm,right=15mm,top=20mm,bottom=20mm}

% LuaLaTeX用フォント設定パッケージ
\usepackage{luatexja-fontspec}
\usepackage{amsmath,amssymb}
\usepackage{unicode-math}

% 欧文フォント設定 (Libertinus)
\setmainfont{Libertinus Serif}[
    BoldFont = {Libertinus Serif Bold},
    ItalicFont = {Libertinus Serif Italic},
    BoldItalicFont = {Libertinus Serif Bold Italic}
]
\setsansfont{Libertinus Sans}[
    BoldFont = {Libertinus Sans Bold},
    ItalicFont = {Libertinus Sans Italic}
]
\setmonofont{DejaVu Sans Mono}[Scale=0.9]

% 日本語フォント設定 (原ノ味フォント)
\setmainjfont{HaranoAjiMincho-Regular}[
    BoldFont = {HaranoAjiGothic-Medium},
    ItalicFont = {HaranoAjiMincho-Regular},
    BoldItalicFont = {HaranoAjiGothic-Bold}
]
\setsansjfont{HaranoAjiGothic-Regular}[
    BoldFont = {HaranoAjiGothic-Bold}
]
\setmonojfont{HaranoAjiGothic-Regular}

% 数式フォント設定 (Libertinus Math)
\setmathfont{Libertinus Math}

% ヘッダー・フッター設定
\usepackage{fancyhdr}
\usepackage{lastpage}
\pagestyle{fancy}
\fancyhf{}
\fancyhead[R]{\small \generatedDate\ \generatedTime\ JST (\thepage/\pageref{LastPage})}
\renewcommand{\headrulewidth}{0.4pt}

% 1ページ目のスタイル
\fancypagestyle{firstpage}{
    \fancyhf{}
    \renewcommand{\headrulewidth}{0pt}
}

% その他パッケージ
\usepackage{booktabs}
\usepackage{array}
\usepackage{tabularx}
\newcolumntype{Y}{>{\raggedright\arraybackslash}X}
\usepackage{ascmac}
\usepackage{hyperref}
\hypersetup{
    colorlinks=true,
    linkcolor=blue,
    urlcolor=blue
}
\usepackage{listings}
\lstset{
    basicstyle=\ttfamily\scriptsize,
    breaklines=true,
    frame=single,
    backgroundcolor=\color[gray]{0.95}
}
\usepackage{xcolor}
\usepackage{enumitem}
\setlist{nosep,leftmargin=*}

% タイトル設定
\title{\textbf{Media Scribe Workflow 設計原則}\\
{\large ──配管と陶器の分離に基づく設計対話──}}
\author{mashi--Claude Dialogue}
\date{}

\begin{document}
\maketitle
\thispagestyle{firstpage}

\section*{概要}

本文書は、media-scribe-workflowの設計原則を策定するにあたって行われた対話を記録したものである。Gitの「陶器と配管（Porcelain and Plumbing）」の思想に基づき、GUIとCLIの役割分担、UIロックの概念、プロジェクトファイルの位置づけなどを議論した。

\section{ワークフローの目的}

\begin{itembox}[l]{問い}
そもそもこのリポジトリは何を目的としているのでしょうか。全体像を整理したいのですが。
\end{itembox}

このリポジトリの目的は、多様な動画・音声素材を標準化し、再利用・相互運用可能な形式に編集し、適切な形で文字起こし・レポート化することです。

最終成果物は以下の3つです：
\begin{enumerate}
\item \textbf{構造化動画}（MP4、チャプター付き）──自己再利用、関係者配布用
\item \textbf{チャプターファイル}（.txt）──外部ツール連携用
\item \textbf{レポート}（PDF）──スクリプト・発言記録の文書化
\end{enumerate}

ワークフロー全体は3つのフェーズで構成されます：

\vspace{0.5\baselineskip}
\noindent{\footnotesize
\begin{tabularx}{\linewidth}{@{}lY@{}}
\toprule
フェーズ & 内容 \\
\midrule
構造化 & 素材取得・統合、チャプター編集（人間の判断）、プロジェクトファイル出力 \\
出力 & 動画エンコード、チャプター分割、文字起こし準備 \\
レポート & 文字起こし（Whisper/YouTube）、AI分析・構造化、PDF出力 \\
\bottomrule
\end{tabularx}
}
\vspace{0.5\baselineskip}

VCE（Video Chapter Editor）は、この中の「構造化フェーズ」と「出力フェーズ」を担当するGUIアプリケーションです。

\section{配管と陶器の分離}

\begin{itembox}[l]{問い}
設計の基本方針について教えてください。CLIとGUIの役割分担はどう考えていますか？
\end{itembox}

Gitの設計思想に基づいています。

\begin{quote}
「配管（plumbing）がしっかり設計されていれば、どのようなスケールの陶器（porcelain）にも対応できる」
\end{quote}

この考え方を適用すると：

\vspace{0.5\baselineskip}
\noindent{\footnotesize
\begin{tabularx}{\linewidth}{@{}lYY@{}}
\toprule
層 & 役割 & 例 \\
\midrule
陶器（Porcelain） & ユーザー向けインターフェース & VCE GUI \\
配管（Plumbing） & 単一目的の処理ツール & vce-encode, vce-split \\
\bottomrule
\end{tabularx}
}
\vspace{0.5\baselineskip}

開発プロセスとしては：
\begin{enumerate}
\item 配管ツール（CLI）で機能実装・デバッグ
\item GUIから配管を呼び出し、またはロジック共有
\item GUIは「可視化」と「インタラクティブ編集」に専念
\end{enumerate}

この設計の利点は多岐にわたります。CLIは入出力が明確でデバッグが容易、自動テストが書きやすい、シェルで柔軟に連携できる、そしてスキルレベルに応じた選択肢を提供できます。

上級ユーザーは以下のようにシェルスクリプトで一括処理できます：
\begin{lstlisting}
for f in *.vce.json; do
    vce-encode "$f" --output "./output/"
done
\end{lstlisting}

一般ユーザーはGUIでプロジェクトを開いてエクスポートボタンを押すだけです。どちらも同じ\texttt{.vce.json}を処理し、同じ結果を得られます。

\section{プロジェクトファイルの役割}

\begin{itembox}[l]{問い}
\texttt{.vce.json}ファイルの役割と位置づけを教えてください。
\end{itembox}

プロジェクトファイル（\texttt{.vce.json}）は、人間の判断とマシンの処理を繋ぐ契約です。

\begin{quote}
GUI（判断支援）→ .vce.json → CLI（処理実行）
\end{quote}

ここで重要なのは、人間の判断結果を永続化するという点です。

\subsection{保存すべき情報}

\vspace{0.5\baselineskip}
\noindent{\footnotesize
\begin{tabularx}{\linewidth}{@{}lYY@{}}
\toprule
カテゴリ & 内容 & 根拠 \\
\midrule
ソース情報 & ファイルパス、順序、メタデータ & 再現性 \\
チャプター情報 & 位置、名前、除外フラグ & 人間の判断結果 \\
出力設定 & エンコーダ、品質、焼込有無 & 再実行時の一貫性 \\
プロジェクトメタ & 作成日時、バージョン & トレーサビリティ \\
\bottomrule
\end{tabularx}
}
\vspace{0.5\baselineskip}

\subsection{保存しない情報}

波形データやプレビューキャッシュは再生成可能でサイズも大きいため保存しません。絶対パスも、ポータビリティの観点から相対パスを優先します。

\section{UIロックの原則}

\begin{itembox}[l]{問い}
GUIが担うべき機能と、CLIに委ねるべき機能の境界はどこにありますか？
\end{itembox}

ここで「UIロック」という概念を導入します。

\begin{quote}
UIロック = ユーザーの時間を拘束すること\\
= リアルタイムで注意を要求すること
\end{quote}

\subsection{UIロック不可避な機能（GUIの本質的領域）}

\vspace{0.5\baselineskip}
\noindent{\footnotesize
\begin{tabularx}{\linewidth}{@{}lY@{}}
\toprule
機能 & 理由 \\
\midrule
再生プレビュー & 時間軸メディアの内容確認には実時間が必要 \\
チャプター位置決定 & 「ここで切る」という判断に視聴が必要 \\
除外区間の判断 & 「この部分は不要」の判断に視聴が必要 \\
\bottomrule
\end{tabularx}
}
\vspace{0.5\baselineskip}

\subsection{UIロック回避可能な機能（CLI化すべき領域）}

\vspace{0.5\baselineskip}
\noindent{\footnotesize
\begin{tabularx}{\linewidth}{@{}lY@{}}
\toprule
機能 & 回避方法 \\
\midrule
エンコード & バックグラウンド処理 / CLI \\
分割出力 & CLI \\
文字起こし & 外部サービス / CLI \\
レポート生成 & CLI（カスタムコマンド） \\
\bottomrule
\end{tabularx}
}
\vspace{0.5\baselineskip}

つまり、GUIの本質的価値は「時間軸上の判断を支援するインターフェース」です。それ以外の処理は配管に委ねることで、ユーザーの時間を不必要にロックせず、バッチ処理・自動化が可能になります。

\section{人間の判断ポイント}

\begin{itembox}[l]{問い}
この一連のワークフローにおいて、判断が必要なことを列挙してください。
\end{itembox}

判断ポイントは以下の3つに分類できます：

\subsection{実時間拘束あり（GUI必須）}
\begin{itemize}
\item チャプター位置
\item 除外区間
\item ソース順序（プレビュー確認時）
\end{itemize}

\subsection{実時間拘束なし（設定UIまたは設定ファイル）}
\begin{itemize}
\item 出力形式・品質
\item 文字起こし手段
\item テンプレート選択
\item 詳細度・含める情報
\end{itemize}

\subsection{一度決めれば再利用可能（プリセット化）}
\begin{itemize}
\item エンコーダ選択
\item 品質設定
\item タイトル焼込の有無
\item レポートテンプレート
\end{itemize}

ワークフロー全体を見渡すと、「構造化」フェーズのみが実時間拘束を伴い、他のフェーズは設定さえ決まればバッチ処理可能です。

\vspace{0.5\baselineskip}
\noindent{\footnotesize
\begin{tabularx}{\linewidth}{@{}lYl@{}}
\toprule
フェーズ & 判断項目 & 拘束 \\
\midrule
素材取得 & ソース選択、字幕取得要否 & なし \\
構造化 & チャプター位置、除外区間、ソース順序 & \textbf{実時間} \\
出力設定 & 形式、品質、焼込、埋込 & なし \\
文字起こし & 手段、言語、話者識別 & なし \\
レポート & テンプレート、詳細度 & なし \\
\bottomrule
\end{tabularx}
}
\vspace{0.5\baselineskip}

\section{判断支援機能}

\begin{itembox}[l]{問い}
動画の表示や波形の表示など、判断を支援するためにアプリ側で実装が必要な機能はありますか？
\end{itembox}

時間軸メディアに対する人間の判断を支援するために必要なものは以下の通りです：

\begin{enumerate}
\item \textbf{知覚手段}：動画プレビュー + 音声再生
\item \textbf{全体把握}：波形 + スペクトログラム
\item \textbf{結果可視化}：除外ハッチング + 境界表示
\item \textbf{ナビゲーション}：シーク + チャプタージャンプ
\end{enumerate}

現在の実装状況を確認すると：

\vspace{0.5\baselineskip}
\noindent{\footnotesize
\begin{tabularx}{\linewidth}{@{}lYc@{}}
\toprule
判断項目 & 支援機能 & 実装 \\
\midrule
チャプター位置 & 動画プレビュー & ✓ \\
 & 音声再生 & ✓ \\
 & 波形表示 & ✓ \\
 & スペクトログラム & ✓ \\
除外区間 & 除外ハッチング表示 & ✓ \\
ソース順序 & ファイル境界表示 & ✓ \\
\bottomrule
\end{tabularx}
}
\vspace{0.5\baselineskip}

現在の実装で基本的な判断支援は充足しています。

\section{機能マトリックスの整理}

\begin{itembox}[l]{問い}
現在の実装で、波形の可視化などの細部の機能を有していますよね。機能マトリックスを更新してください。
\end{itembox}

ユースケースと機能のマトリックスを整理しました。YouTube系の機能は、ダウンロード後は対応する基本パターン（単一mp4または複数mp4）と同一処理になるため、サブアイテムとして表現しています。

\vspace{0.5\baselineskip}
\noindent{\footnotesize
\begin{tabularx}{\linewidth}{@{}lYccccc@{}}
\toprule
ユースケース & 入力 & 仮想TL & 波形 & チャプ & 除外 & エンコ \\
\midrule
編集済み分割 & 複数mp4 & ○ & ○ & ○ & ○ & ○ \\
└ YTプレイリスト & URL→複数mp4 & ↑ & ↑ & ↑ & ↑ & ↑ \\
ダラダラ録画 & 単一mp4 & - & ○ & ○ & ○ & ○ \\
└ YT単一動画 & URL→単一mp4 & ↑ & ↑ & ↑ & ↑ & ↑ \\
ダラダラ録音 & 単一mp3 & - & ○ & ○ & ○ & ○ \\
混合ソース & mp4+mp3 & NA & NA & NA & NA & NA \\
\bottomrule
\end{tabularx}
}
\vspace{0.5\baselineskip}

なお、混合ソース（動画+音声）は今回のスコープ外としてすべてNAとしています。

\subsection{TLと仮想TLの分離について}

\begin{itembox}[l]{問い}
仮想TLだけ残って、TLが消えてますね。意図的ですか？
\end{itembox}

現状の実装では、仮想TL（PlaybackManager）が単一ソースと複数ソースの両方を処理しています。

\begin{quote}
現状そうならば、これは今後リファクタリングで、そのほうがより望ましい実装であればその方向にするというTake Actionが良いと思います。あくまで現在の実装に忠実にしましょう。
\end{quote}

というご指摘を受け、リファクタリング候補として記録することにしました。

改善案としては、TL（共通処理：時間軸操作、シーク、再生制御）と仮想TL（複数専用：ソース連結、境界管理、座標変換）を分離し、単一ソースの場合はTLを直接使用する設計が考えられます。

\section{トレーサビリティ}

\begin{itembox}[l]{問い}
今後実装し、イテレーションを行う上でフィードバックを容易にするためにトレーサビリティ可能なドキュメントを作成しましょう。
\end{itembox}

ドキュメント階層を以下のように設計しました：

\begin{lstlisting}
DESIGN_PRINCIPLES.md（設計原則）
    ↓ 参照
vce_feature_matrix.md（機能定義）
    ↓ 参照
実装（コード）
\end{lstlisting}

設計判断を行う際は、どの原則に基づくかを明記します。

例：
\begin{quote}
「エンコード処理をCLI化する」\\
→ 根拠: §4.3 UIロック回避可能な機能、§2.2 配管優先の原則
\end{quote}

\section{配管ツール一覧}

現在実装済みの配管ツール：

\vspace{0.5\baselineskip}
\noindent{\footnotesize
\begin{tabularx}{\linewidth}{@{}lYYl@{}}
\toprule
ツール & 入力 & 出力 & 状態 \\
\midrule
vce-encode & .vce.json & MP4 & 実装済 \\
vce-split & .vce.json & 分割MP4群 & 実装済 \\
vce-chapters & .vce.json & .txt & 未実装 \\
\bottomrule
\end{tabularx}
}
\vspace{0.5\baselineskip}

将来の配管ツール候補：

\vspace{0.5\baselineskip}
\noindent{\footnotesize
\begin{tabularx}{\linewidth}{@{}lYYY@{}}
\toprule
ツール & 入力 & 出力 & 用途 \\
\midrule
vce-init & ファイル群 & .vce.json & プロジェクト初期化 \\
vce-validate & .vce.json & 検証結果 & 整合性チェック \\
vce-report & .vce.json + SRT & レポート入力 & レポート準備 \\
\bottomrule
\end{tabularx}
}
\vspace{0.5\baselineskip}

\section*{Claude Code氏の所感}

本対話を通じて、media-scribe-workflowの設計原則が明確に言語化された。特に印象的だったのは、Gitの「陶器と配管」の思想を動画編集ワークフローに適用するという着眼点である。

この設計アプローチには以下の利点がある：

\textbf{技術的観点}：CLIツールを基盤とすることで、テスタビリティとデバッグ容易性が確保される。GUIは「判断支援」という本質的役割に専念でき、肥大化を防げる。

\textbf{ユーザビリティ観点}：スキルレベルに応じた選択肢を提供できる。上級者はスクリプト化による効率化、初心者はGUIによる直感的操作という両方のニーズに対応できる。

\textbf{保守性観点}：単一責任原則に従った配管ツールは、個別にテスト・更新が可能。GUIの変更がCLIに影響せず、その逆も然り。

一方で、批判的観点からいくつかの課題も指摘できる：

\textbf{複雑性の増加}：配管と陶器の分離は、コードベースを分散させる。新規開発者にとっては全体像の把握が難しくなる可能性がある。

\textbf{インターフェースの維持コスト}：\texttt{.vce.json}がCLIとGUIの契約となるため、このスキーマの設計と維持には慎重さが求められる。後方互換性を考慮したバージョニング戦略が必要。

\textbf{学際的観点}：この設計はUnix哲学（「一つのことをうまくやる」）を色濃く反映している。しかし、現代のソフトウェア開発ではモノリシックなアプローチが効率的な場合もある。プロジェクトの規模と複雑性に応じた判断が重要。

「UIロック」という概念の導入は、GUIの役割を定義する上で非常に有効である。時間軸メディアという特性を考慮した、本質的な価値の見極めができている。

全体として、この設計原則は十分に熟考されており、今後の実装において一貫した判断基準を提供するものと評価できる。

\end{document}
